\section{Dataset and Data Processing}
\label{sec:methodology}

Data processing is critical in FER, particularly under the low input resolution of 64x64 pixels. At this resolution, fine-grained texture information is limited. While benchmarks like FER2013 successfully use even smaller 48x48 grayscale images \cite{khaireddin2021facialemotionrecognitionstate}, recent studies suggest that higher resolutions (e.g., 256x256) can improve generalization \cite{tutuianu2023benchmarking}. Our preprocessing must prioritize precise face detection and alignment to maximize information density in the available pixels. We intend to use automated face detection followed by landmark-based alignment (e.g., eye-centering via an affine transformation) to ensure that key facial regions such as the eyes and mouth are consistently positioned across samples. To maintain consistency across datasets, all images are resized to a fixed resolution and converted to a unified color format (RGB), resulting in a standardized input representation.

To tackle overfitting, a risk with limited datasets, we plan to use a mix of data augmentation techniques. Standard geometric transformations are essential. Specifically, we may apply rotation ($\pm$ 10°) and horizontal flipping, as suggested by Khaireddin \& Chen \cite{khaireddin2021facialemotionrecognitionstate}. Vertical flipping is excluded as it contradicts the natural orientation of faces. Beyond geometry, we will implement pixel-level perturbations like random erasing \cite{ramos-cosi2025realtime}. We also consider masking techniques, which recent surveys identify as strong self-supervised learning signals that force the network to learn context rather than relying on local artifacts \cite{li2023deep}. Class imbalance is also a challenge in FER: Often, expressions such as \textit{Happiness} are overrepresented, while \textit{Disgust} and \textit{Fear} are more rare. Naive inverse-frequency weighting can overemphasize minority classes and lead to overfitting on noisy samples. Therefore, we adopt \textit{Conservative Class Weighting}, which caps class weights at a predefined threshold to avoid exaggerated gradients while still compensating for imbalance \cite{mejiaescobar2023towards}. We also consider combining multiple datasets to reduce imbalance and increase diversity. When combining datasets, we align label definitions and class distributions to reduce domain shift. We consider publicly available (sub-)sets of established benchmarks such as AffectNet, RAF-DB, and FER2013, with AffectNet in particular delivering greater in-the-wild variety. Finally, we recognize established demographic biases in FER datasets \cite{gayamorey2025evaluating, ghotbi2022ethics}. To uncover possible systematic differences, model performance is assessed globally and across available demographic subsets.

\section{Conclusion and Outlook} 
% \label{sec:conclusion}

This initial report inspired appropriate data processing strategies and model architectures for recognizing low-resolution facial expressions. Next, we plan to create a solid convolutional baseline with a ResNet-based architecture to confirm our choices for preprocessing and augmentation. We will assess attention-augmented CNNs, as well as transformer-based or hybrid models, using this baseline for evaluation. In addition to overall accuracy, performance will be evaluated using class-wise metrics and fairness-oriented analyses. The model with the best performance will be refined for real-time inference and incorporated into a demonstration using a webcam.